% =============================================================================
%  The Construction of the FTD Mathematics
%  From a Discrete Ontology to Its Provable Boundary
% -----------------------------------------------------------------------------
%  Master file.  Build (Perl-free; latexmk unavailable on this MiKTeX):
%     pdflatex monograph  ->  biber monograph  ->  pdflatex monograph  x2
%  or run  build.ps1 / build.bat  in this directory.
%
%  This file is a standalone typeset monograph.  It introduces no new
%  mathematics and promotes no epistemic tag.
% =============================================================================
\documentclass{ftdmonograph}

\usepackage{ftd-math-macros}     % canonical equations & constants (anti-drift)
\usepackage{ftd-epistemic}       % the monochrome epistemic-tag system

\addbibresource{references.bib}

\begin{document}

\frontmatter
% =============================================================================
% frontmatter.tex  --  title page, scope/authorities, how-to-read + rule key, ToC
% =============================================================================

% ---- title page -----------------------------------------------------------
\begin{titlingpage}
  \centering
  \vspace*{\stretch{1.2}}
  {\Huge\scshape The Construction of\\[3pt] the FTD Mathematics\par}
  \vspace{14pt}
  {\large\itshape From a Discrete Ontology to Its Provable Boundary\par}
  \vspace{20pt}
  {\rule{0.46\textwidth}{0.4pt}\par}
  \vspace{20pt}
  {\small\scshape a monograph on foundational ternary dynamics\par}
  \vspace{\stretch{1}}
  \begin{minipage}{0.82\textwidth}
    \centering\itshape\small
    Derive everything the discrete ontology forces --- and rigorously
    establish the boundary of what it cannot force.
  \end{minipage}\par
  \vspace{\stretch{1.4}}
  % NOTE (O-1): author line deliberately omitted. The source monograph carries
  % no author; the project's grade-A paper leaves \author{} empty for the owner
  % to fill at submission time. Do NOT fabricate a name.
  {\small 2 June 2026\par}
  \vspace*{\stretch{0.4}}
\end{titlingpage}

% ---- scope ----------------------------------------------------------------
\section*{Scope}
\noindent This monograph builds FTD's mathematics from the ground up, in the
order a mathematician would build the number systems, for an interdisciplinary
and philosophy-of-physics readership. It is a construction and an honest
accounting at once: it introduces no new mathematics, and it states every result
at exactly the grade the result has earned --- theorem, derivation, conjecture,
or open problem --- and at no grade higher. That discipline is not a hedge; it is
the point. It lets the strong results stand as strong and the boundary stand as a
genuine finding.

% ---- how to read ----------------------------------------------------------
\section*{How to read this monograph}
\noindent This is a \emph{construction story}. It builds FTD's mathematics in
the order a mathematician would build the number systems --- each object from
prior ones, each step independently checkable --- and it asks one question at
every step: \textbf{does the discrete ontology force this, or not?} (where
``the discrete ontology'' includes the one $\mathbb{Z}[i]$ modelling choice of
\secref{sec:seed}). The answer comes with an epistemic tag, and the tag is the
point. A reader from philosophy of physics is invited to scrutinize the tags
first and the prose second; the tags are where the project's honesty lives.

\medskip
\noindent The monograph is organized around the project's single governing
goal, stated as two clauses:
\begin{quote}\itshape
  Derive everything the discrete ontology forces --- and rigorously establish
  the boundary of what it cannot force.
\end{quote}
\noindent Both clauses are deliverables. Showing that the discrete ontology
\emph{does not} fix a quantity is as much a result of this program as showing
that it does; a mapped boundary is not a failure but a finding.

\medskip
\noindent\textbf{How a tag looks.} Throughout, an epistemic tag is set as a
small-caps badge whose \emph{rule} --- never its colour --- carries its tier,
so the page reads identically in grayscale. There are four treatments:
\rulekeylegend
\noindent The same rule signatures frame the displayed ``marquee'' statements
and label the construction map of \cref{fig:dag}. The full contract --- every
tag, its meaning, and what a reviewer should expect of it --- is the table of
\secref{sec:contract}.

\cleardoublepage
\tableofcontents*


\mainmatter
\setcounter{part}{-1}            % so the first \part renders as "Part 0"

\part{The Seed}% =============================================================================
% Part 0 -- The Seed
% =============================================================================
\noindent Part~0 lays the foundation the rest of the monograph stands on: the
discrete ontology FTD posits, the standard of ``construction'' the document
holds itself to, the tag system the reader must internalize, and the two-clause
roadmap that organizes everything that follows. It contains \textbf{no physics
theorem and no new mathematics}; it is setup.

\section{The discrete ontology}\label{sec:ontology}
FTD posits a discrete substrate, fixed by five postulates. Each is an
\etag{ax}: a structural commitment that defines the model and is
\textbf{not} claimed to be derived from anything more primitive. A reader
should accept them as the definition of the game, then ask what follows.
\margetag{ax}%
\begin{enumerate}[leftmargin=2.2em, itemsep=3pt, label=\arabic*.]
  \item \textbf{Discrete space.} The substrate is a three-dimensional cubic
        lattice with \textbf{undefined boundary}: at every specified position,
        the axis-adjacent sites exist. This is deliberately \emph{not} a
        completed-infinity totality $\mathbb{Z}^3$ taken as a single finished
        object; it is a lattice with no defined edge and no claim of
        completeness.
  \item \textbf{Discrete time.} Evolution proceeds in discrete ticks; there is
        no continuous time parameter at the substrate level.
  \item \textbf{Ternary states.} Each site holds a state $s\in\{-1,0,+1\}$, the
        three values being the only states a site may take.
  \item \textbf{Local causality.} Influence is local: each site interacts only
        within its 26-neighbour Moore neighbourhood, and information propagates
        at most one lattice unit per tick.
  \item \textbf{Determinism.} The update rule is deterministic: the
        configuration at tick $t+1$ is fixed by the configuration at tick $t$.
\end{enumerate}

\runin{The two-layer ontology \etag{ax}.} Riding on the lattice are two fields
with distinct ontological roles:
\begin{itemize}[leftmargin=1.4em, itemsep=2pt]
  \item a \textbf{flux field} $J\in\R^3$, a continuous vector field at each
        site, encoding potential energy density. It is \emph{dispositional}: it
        represents what the substrate is poised to do, not what it has manifested.
  \item a \textbf{state field} $s\in\{-1,0,+1\}$, the discrete ternary value
        above. It is \emph{actual}: it represents manifestation, what the
        substrate has resolved to.
\end{itemize}
The dispositional/actual distinction is structural to the framework and recurs
throughout; the flux field carries the continuous content, the state field
carries the discrete resolution.

\runin{Why the undefined boundary is load-bearing.} The choice of an
undefined-boundary lattice over a completed-infinity $\mathbb{Z}^3$ is not
cosmetic; it narrows the framework's admissible toolkit and is enforced
corpus-wide. Concretely, statements of the form ``in the limit $L\to\infty$
\ldots'' are \textbf{not well-posed} under this ontology without an explicit
$\varepsilon$--$L$ restatement (e.g.\ ``for every $\varepsilon>0$ there exists
$L$ such that $P(L)$ is within $\varepsilon$ of its continuum counterpart'').
Global integrals over ``all sites,'' completed thermodynamic limits, and path
integrals over ``all configurations'' are likewise not primitive moves;
arbitrarily large \emph{finite} computations, explicitly bracketed, are. This
commitment will matter precisely at the boundary in Part~II, where the question
``does the ontology force a particular limiting value?'' turns on what limits the
ontology is even entitled to take.

\section{What a construction is}\label{sec:construction}
This monograph holds itself to a specific standard, and naming it up front is
part of the contract with the reader.

\runin{A construction,} in the sense used here, is the kind of thing the number
systems are: $\N\to\Z\to\Q\to\R\to\C$. Each object is built \emph{explicitly}
from objects already in hand ($\Z$ as equivalence classes of pairs of naturals;
$\Q$ as equivalence classes of pairs of integers; $\R$ as cuts or Cauchy
sequences; $\C$ as $\R^2$), and each construction step is \emph{independently
checkable} by a reader who accepts the prior stage. No object appears by fiat;
none is assumed because it would be convenient. FTD's mathematical core is
presented in exactly this spirit: each result built from prior results and
the axioms, with a stated proof or chain the reader can audit. The ordering is
fixed: mathematical primitives $\to$ invariant structure $\to$ admissible
readouts $\to$ operational physics, with physical language entering \emph{only}
at the last layer. The constructive Parts of this monograph (I and II) live in
the first two layers; the physics \emph{match} is withheld until Part~III, and
Part~II names the physical target only to state precisely what is \emph{not}
forced.

\medskip
\noindent Three categories must be kept strictly apart, because conflating them
is the characteristic way a program of this kind overclaims:
\begin{itemize}[leftmargin=1.4em, itemsep=4pt]
  \item \textbf{A derivation} (tagged \etag{thm} or \etag{der}) is an
        \emph{explicit chain} from the axioms or prior theorems that the
        document itself reproduces. The reader checks the chain. \etag{der} is
        the honest grade when the chain is explicit but carries a non-trivial
        assumption; \etag{thm} when it does not.
  \item \textbf{A parametric insertion} (tagged \etag{par}) is a
        standard physics formula with FTD's constants substituted in. The
        numbers may fit beautifully, but the \emph{mechanism is borrowed} from
        outside the framework: a calibration input, not a derivation.
  \item \textbf{A match} (at its strongest tagged \etag{smc}) is a
        numerical coincidence between an FTD object and a physical quantity,
        possibly backed by uniqueness scans or sub-ppm agreement, but with
        \textbf{no derivation chain}. A match is evidence, not a derivation.
\end{itemize}
For every load-bearing claim, this monograph states which of the three it is.

\section{The epistemic contract}\label{sec:contract}
The following tag system is the instrument the whole program is governed by. A
reader should internalize it before proceeding; every substantive claim
sentence in this monograph carries one of these tags, and the tag (not the
surrounding prose) is the claim's actual epistemic status.

\begingroup
\footnotesize
\renewcommand{\arraystretch}{1.25}
\setlength{\LTpre}{8pt}\setlength{\LTpost}{8pt}
\begin{xltabular}{\linewidth}{@{}l >{\RaggedRight\arraybackslash}X >{\RaggedRight\arraybackslash}X@{}}
  \toprule
  \textsc{tag} & \textsc{meaning} & \textsc{reviewer expectation}\\
  \midrule
  \endfirsthead
  \toprule
  \textsc{tag} & \textsc{meaning} & \textsc{reviewer expectation}\\
  \midrule
  \endhead
  \bottomrule
  \endlastfoot
  \ctag{ax} & Structural postulate; not derivable. & Accept as model definition.\\
  \ctag{thm} & Rigorously proven from axioms (or from named classical theorems
    with explicit citation). & Check the proof.\\
  \ctag{der} & Established by an explicit chain the document reproduces; weaker
    than a theorem when the chain carries a non-trivial assumption. & Check
    the chain; flag any smuggled axiom.\\
  \ctag{sel} & Argued from consistency or naturalness; not uniquely proven.
    & Critique the argument.\\
  \ctag{smc} & A conjecture with substantial structural and/or empirical
    evidence (uniqueness scan, multi-route convergence, sub-ppm match) but
    \textbf{no} derivation chain. & Critique the evidence; expect an explicit
    Bayes / uniqueness / look-elsewhere argument.\\
  \ctag{cnj} & Proposed interpretation requiring validation; weaker than the
    above (no structural-uniqueness backing). & Demand evidence.\\
  \ctag{par} & A standard physics formula filled with FTD constants; numbers
    fit, mechanism borrowed. & Treat as calibration input, not output.\\
  \ctag{cn} & A hypothesis tested and falsified; preserved for provenance to
    prevent re-attempt. & Confirm the closure; cite it to prevent a zombie
    re-emergence.\\
  \ctag{syn} & Cross-document integration of existing claims at their existing
    tags; not a new theorem. & Verify the component claims; check that nothing
    was silently promoted.\\
\end{xltabular}
\endgroup

\runin{The reading rule.} The tags form a ladder, and the ladder has a
direction:
\begin{quote}\itshape
  Promotion up the ladder requires explicit justification. Demotion is free.
  Ambiguous cases default down.
\end{quote}
\noindent A claim may be moved \emph{up} (e.g.\ from \etag{cnj} to \etag{der})
only by exhibiting the chain or proof that warrants it. A claim may be moved
\emph{down} at any time, by anyone, for any defensible reason, with no ceremony.
And where the correct tag is genuinely unclear, the convention is to assign the
\emph{weaker} tag until the stronger one is earned. This asymmetry is the
safeguard against the program talking itself into believing its own conjectures.

\medskip
\noindent\textbf{Every result appears at the grade it has earned, and at no
grade higher.} (Four further tags appear in the body where they are
needed: \etag{meas} for an honest measurement, \etag{imp} for an imposed
choice, \etag{opn} for an open problem, and \etag{fo} for a foundational
obstruction.)

\section{The two clauses: roadmap}\label{sec:roadmap}
The program's governing goal has two clauses, and the architecture of this
monograph is simply those two clauses made into sections, plus a quarantined
bridge to physics.
\begin{quote}\itshape
  Clause 1: Derive everything the discrete ontology forces.\\
  Clause 2: Rigorously establish the boundary of what it cannot force.
\end{quote}

\runin{Clause 1 $\to$ Part I, the Constructive Reach.} Part~I exhibits what the
five postulates \emph{force} once their order-4 lattice symmetry is read as
$\Zi$ (the single modelling \etag{ax} of \secref{sec:seed}): an algebraic spine
of theorem-grade results: the constant $\Gs=\Gamma(1/4)/\Gamma(3/4)$, the master quadratic
$\masterquad$ and its roots, the arithmetic of the lemniscatic CM curve, and the
finite combinatorics of the Moore neighbourhood. The canonical accounting is
\textbf{nine numbered results, of which seven are theorem-grade} and two are
honestly tiered below theorem grade. These are the framework's firm ground,
independent of any physics interpretation, and Part~I states each at its
canonical tag.

\runin{Clause 2 $\to$ Part II, the Boundary, and this is the climax.}
Part~II asks the sharp negative question: granted everything in Part~I, does the
discrete ontology \emph{force} the electromagnetic coupling $\alpha$? The honest
answer, mapped across multiple independent attack routes, is \textbf{no, not
without an additional binding law.} The boundary takes a concrete and memorable
form: the reading that would fix $\alpha$ requires the substrate to
\emph{natively realize} the quantity $\sqrtwall$ (the discrete limit of a
square root) as an operator-assembly the lattice produces on its own, and
across every FTD-native route examined the lattice does not produce it; the
determinant grading and the operator assembly $\assembly$ are an imposed
selection, not a forced consequence. The boundary itself is built from
theorem-grade parts around one precisely-stated open kernel: a ruled-out geometric realization, a closed-scope
restriction, and a route-invariant reduction (the same kernel reached in four
different dialects, because $\Q(\Gs)$ is the Galois-fixed field of the master
quadratic's $\Z/2$). Part~II's
deliverable is therefore a \emph{mapped wall}: theorem-grade masonry around one
clearly-marked gap, exactly the kind of boundary Clause~2 calls for.

\runin{The bridge $\to$ Part III, quarantined.} The numerical match that
motivates the whole enterprise ($x_+ = \xplusnum$, agreeing with
$1/\alpha$ to $\ppmval$~ppm) is real and is the strongest structural
evidence the framework holds (a unique dual-matcher across millions of candidate
polynomials over a basket FTD did not design). But it is a \emph{match}: it
carries the tag \etag{smc}, not a derivation, and Part~III presents it as
such, alongside the \etag{par} physics insertions.

\runin{The central honest headline, stated up front.} The discrete ontology,
\textbf{once its order-4 lattice symmetry is read as $\Zi$ (the single modelling
\etag{ax} isolated in \secref{sec:seed})}, forces a rich algebraic
spine (\textbf{seven theorem-grade results}), and Part~II proves the precise
conditional boundary around the electromagnetic coupling $\alpha$: the present
theory does not select it without an additional binding law. Both halves are the
deliverable. A program that only proved the
first half would be a curiosity; a program that mapped the second half honestly
is doing the harder and more useful thing: drawing, in both directions, the
line of what discreteness reaches. The constant $\Gs\approx\Gstarapprox$ that
anchors the spine is, throughout, the lemniscatic ratio $\Gamma(1/4)/\Gamma(3/4)$;
it is \textbf{not} the Bernoulli/Gauss lemniscate constant
$\varpi\approx\varpiapprox$, a distinct number with which it is never to be
conflated.

\part{The Constructive Reach}% =============================================================================
% Part I -- The Constructive Reach
% =============================================================================
\noindent Part~I answers Clause~1: \textbf{what does the discrete ontology
force?} It exhibits the algebraic spine (the theorem-grade results the
substrate produces) in the order a mathematician would build them, each
object from prior ones, each step independently checkable. The construction runs
$i \to \Zi \to \Gs \to \text{master quadratic} \to \text{spine}$. It carries
\textbf{no physics identification}: the famous numerical match $x_+ = 1/\alpha$
is a \emph{match}, deferred to Part~III, and Part~I says so at the one place it
is tempting to overstep.

\medskip
\noindent The spine is \textbf{nine numbered results: seven are theorem-grade,
and two are honestly tiered below theorem grade} (the coefficient-16 identity,
whose structural necessity is conjectural, and the Phase-J ultralocality, which
is a theorem only at $L=2$). \Cref{sec:spine} closes by re-stating that count.

\section{The seed: \texorpdfstring{$i$}{i} and \texorpdfstring{$\Zi$}{Z[i]}}\label{sec:seed}
The construction begins with one structural choice and one group-theory fact,
and it is worth separating them cleanly, because the honesty of everything
downstream depends on not confusing the two.

\runin{The structural choice.}\margetag{ax} The substrate is the cubic lattice of
Part~0. Its coordinate planes carry an order-4 rotational symmetry: a $90^\circ$
rotation of the $(x,y)$ plane sends
$(1,0)\mapsto(0,1)\mapsto(-1,0)\mapsto(0,-1)\mapsto(1,0)$, a cycle of length~4.
FTD's seed move is to \emph{read this order-4 planar symmetry as the arithmetic
of the Gaussian integers} $\Zi=\{a+bi : a,b\in\Z\}$, with the $90^\circ$ rotation
realized as multiplication by $i$. That the lattice's quarter-turn \emph{should
be read through} $\Zi$ rather than through some other order-4 structure is a
modelling decision, not a theorem, the discrete-ontology analogue of choosing
to build $\C$ as $\R[i]$: natural, motivated, and a \emph{choice}. We mark it an
\etag{ax} and flag it as such, so that no later theorem silently inherits the
status of ``forced'' from a step that was actually ``chosen.''

\runin{The group-theory facts.}\margetag{thm} Once $\Zi$ is the substrate's
arithmetic, the following are ordinary, checkable theorems of algebra, owing
nothing further to FTD:
\begin{itemize}[leftmargin=1.4em, itemsep=3pt]
  \item The unit group of $\Zi$ is $\Ziu=\{1,i,-1,-i\}\cong\Z/4$, of order
        $|\Ziu|=4$. \emph{(Proof: $a+bi$ is a unit iff its norm $a^2+b^2=1$,
        whose only integer solutions are $(\pm1,0)$ and $(0,\pm1)$.)}
  \item $\Zi$ is a Euclidean domain, hence a PID and a UFD; its primes split
        into the three Gauss classes: the ramified prime $1+i$ (norm~2), the
        split rational primes $p\equiv 1\ (\mathrm{mod}\ 4)$, and the inert
        rational primes $p\equiv 3\ (\mathrm{mod}\ 4)$ (Fermat's two-square
        theorem).
\end{itemize}
The single integer this seed deposits into the rest of the construction is
$4=|\Ziu|$. It reappears, squared, as the coefficient~16 of the master quadratic
(\secref{sec:masterquad}), and the split/inert residue classes mod~4 reappear as
the arithmetic content of one of the four $\Gs$ constructions (\secref{sec:gstar}).
Hold the distinction firmly: $4$ is a \emph{theorem} about $\Zi$; that the
\emph{lattice} is entitled to $\Zi$ in the first place is the axiom.

\section{\texorpdfstring{$\Gs$}{G*} in four representations of differing provenance}\label{sec:gstar}
The bridge constant is\margetag{thm}
\[ \Gstardef. \]
What earns $\Gs$ its place at the centre of the spine is that it admits four
\emph{representations} of sharply differing character: a $\Gamma$-ratio, a
lattice Green's-function period, a regularized determinant ratio, and a finite
product. These are not four independent witnesses: three share the quarter-integer
data, and the Watson value is itself a $\Gamma(1/4)$-period, so all four are
representations of one and the same CM period of the lemniscatic curve. What is
striking is that a single period wears four such different faces (arithmetic,
lattice-analytic, spectral, and finitary), one of which (Watson's
body-centred-cubic integral) is of genuinely independent provenance. Each is shown
where it is cheap and cited where it is deep.

\runin{Route 1: the $\Gamma$-ratio.} Directly from the definition and the
Euler reflection formula,
\[ \reflection, \]
multiplying $\Gs=\Gamma(1/4)/\Gamma(3/4)$ by $\Gamma(3/4)^2$ and substituting
gives the two standard closed forms
\[ \Gstarclosed, \]
where $\varpidef$ is the Bernoulli/Gauss lemniscate constant. A three-line proof,
verified to 50 digits.

\runin{Route 2: the Watson BCC period}, conditional on Watson 1939. Let $W_3$
be the body-centred-cubic Watson integral
\[ \watsonint. \]
Watson's 1939 closed-form evaluation~\cite{watson1939}, consolidated by
Glasser--Zucker~\cite{glasserZucker1980}, gives
\[ \watsonval. \]
The BCC sub-lattice is one of the three polyhedral layers of the substrate's
Moore neighbourhood, so this is $\Gs$ arising from a \emph{lattice}
quantity, a Green's-function period of an entirely different character from
Route~1's $\Gamma$-ratio. ``Conditional on Watson 1939'' means the rigour equals
that of the cited classical evaluation, not that it rests on a conjecture.
Verified to 100 digits in PARI.

\runin{Route 3: the $\det_\zeta$ quarter-conjugacy bridge.} Let $J$ be the
conjugacy operator with $J^2=-I$ (multiplication by $i$; the same $Z_4$ seed of
\secref{sec:seed}). A wavefunction on $S^1$ with the quarter-twisted boundary
condition $\psi(\varphi+2\pi)=J\,\psi(\varphi)$, evolved by the natural
quarter-period operator $-i\,d/d\varphi$, carries its two $J$-eigensectors on the
shifted spectra $D_{1/4}=\{n+1/4\}_{n\ge0}$ and $D_{3/4}=\{n+3/4\}_{n\ge0}$ (the
choice of that operator is a modelling choice, like the seed, not something
forced). Lerch's formula for the $\zeta$-regularized
determinant of an arithmetic progression,
\[ \lerch, \]
yields the determinant ratio
\[ \detzratio, \]
the $\sqrt{2\pi}$ cancelling. The arithmetic is the reason this route belongs to
the same story as the seed: $4\cdot D_{1/4}=\{n\equiv1\ (\mathrm{mod}\ 4)\}$ and
$4\cdot D_{3/4}=\{n\equiv3\ (\mathrm{mod}\ 4)\}$ are exactly the two non-trivial
residue classes mod~4, and, restricted to primes, the split and inert prime
classes of $\Zi$ (\secref{sec:seed}). $\Gs$ is the regularized asymmetry between
them.

\runin{Route 4: the finite-$N$ attractor.} Define the finite product
\[ \finiteN. \]
Every $G^{*}_N$ is computable in $O(N)$ rational operations with no transcendental
ever invoked ($G^{*}_0=3$, $G^{*}_1=21/(5\sqrt2)\approx2.9698$,
$G^{*}_5\approx2.95995$, $G^{*}_{20}\approx2.95878$). A Stirling expansion of
$\Gamma(N+7/4)/\Gamma(N+5/4)\sim(N+1)^{1/2}$ gives
\[ \finiteNrate \]
(empirical $C\approx0.046$). This route does double duty. It is a fourth
construction of $\Gs$, and it \textbf{discharges the Part-0
$\varepsilon$--$L$ obligation} for $\Gs$: the undefined-boundary ontology forbids
treating $\Gs=\lim_{N\to\infty}G^{*}_N$ as a primitive completed limit, and
Route~4 supplies exactly the finitary restatement the ontology demands: $G^{*}_N$
is defined for every finite $N$ and approaches $\Gs$ at the stated rate. $\Gs$ is
therefore reachable by a convergent finite computation, no completed infinity
required.

\runin{$\Gs$ is the lemniscatic ratio, not $\varpi$.} $\Gs\approx\Gstarapprox$ is
the ratio $\Gamma(1/4)/\Gamma(3/4)$. It is \textbf{not} the Bernoulli/Gauss
lemniscate constant $\varpi\approx\varpiapprox$, a distinct number related by
$\Gs=2\varpi/\sqrt\pi$. Informal usage sometimes calls both ``the lemniscate
constant''; here they are never conflated, and the distinction is load-bearing in
the most literal sense. The master quadratic of \secref{sec:masterquad} produces
$x_+=\xplusnum$ \textbf{only} at $\Gs=\Gstarapprox$; substituting
$\varpi=\varpiapprox$ into the same polynomial gives $x_+\approx107.3$, nowhere
near $1/\alpha$.

\section{The master quadratic}\label{sec:masterquad}
With $\Gs$ in hand, the construction's central polynomial is purely algebraic.

\begin{forcedbox}[The master quadratic \etag{thm}]
Define
\[ \masterquad. \]
Its discriminant factors cleanly,
\[ \discriminant \]
which is \textbf{positive} because $\Gs>1/4$ (indeed $\Gs\approx2.96$), so both
roots are real:
\[ \rootsexplicit \]
Numerically $x_+=\xplusnum$ and $x_-=\xminusnum$.
\end{forcedbox}

\noindent This is the quadratic formula applied to a polynomial whose
coefficients are integer multiples of powers of $\Gs$; nothing more is involved.
Verified numerically across 54 independent checks.

\runin{The coefficient-16 soft spot, flagged honestly.}\margetag{cnj} The
prefactor~16 has a tempting arithmetic reading: for the lemniscatic CM elliptic
curve $E:y^2=x^3-x$, the automorphism group over $\overline{\Q}$ is
$\Aut(E)=\{\pm1,\pm i\}\cong\Z/4$ (the unit group of $\Zi$ again), so
$|\Aut(E)|=4$ and $|\Aut(E)|^2=16$. That $|\Aut(E)|^2=16$ is a \etag{thm}. But
that the master quadratic's prefactor is \emph{forced} to equal $|\Aut(E)|^2$ (that
two objects equal 16 \emph{for a structural reason}) is not proven: it is a
\etag{cnj} on the structural necessity, the framework's softest mathematical
spot. A partial unification exists (a tower-level reading supplies a reason for
$k=4$, hence $2^k=16$, via $|\Ziu|^2=16$), but a partial unification is not a
forcing theorem, and the honest grade stays \etag{cnj} on necessity.

\runin{No physics here.} $x_+=137.036\ldots$ is, \emph{in Part~I}, a root of a
polynomial, nothing more. The proximity to $1/\alpha$ is the motivation for
the whole enterprise, but the identification $x_+ \leftrightarrow 1/\alpha$ is a
\etag{smc} deferred in full to \textbf{Part~III}; it is not asserted as forced
anywhere in the constructive Parts. The smaller root $x_-\approx\xminusapprox$ is
a pure algebraic artifact, pinned by the Vieta relation $x_-=16\Gsss/x_+$ the
instant $x_+$ is, and it carries no physics. The historical identification
$x_-\leftrightarrow N_c$ is \textbf{retired}; FTD's $N_c=3$ is independently
sourced from the Moore-neighbourhood topology, owing nothing to the polynomial's
smaller root.

\section{The spine built on top}\label{sec:spine}
On the foundation of $\Gs$ (\secref{sec:gstar}) and the master quadratic
(\secref{sec:masterquad}), a further family of theorem-grade results stands. Each
is stated at its exact grade and condition; short proofs are shown, deep external
dependencies are cited.

\runin{Harmonic-invariant tower.}\margetag{thm} For each integer $k\ge3$, the
$(1+i)$-tower master quadratic is $\towerquad$ (the $k=4$ instance is
\secref{sec:masterquad}'s polynomial). Writing the normalized roots
$y_\pm := x_\pm/\Gs$, at \emph{every} level $k\ge3$,
\[ \harmonicinv. \]
The proof is a three-line Vieta computation:
$1/x_+ + 1/x_- = (x_+ + x_-)/(x_+ x_-)
= (2^k\Gpow{k-2})/(2^k\Gpow{k-1}) = 1/\Gs$; multiplying through by $\Gs$ gives
$1/y_+ + 1/y_- = 1$. Verified to 50 digits.

\runin{CM-curve uniqueness}, under the trivial-multiplier criterion.\margetag{thm}
Among all imaginary quadratic fields $K=\Q(\sqrt{-d})$, the field $\Q(i)$ ($d=1$,
fundamental discriminant $-4$) is the \textbf{unique} one satisfying
$|\mu_K|=|\mathrm{disc}(K)|$: for $\Q(i)$, $|\mu_K|=4$ and
$|\mathrm{disc}(K)|=4$; for every other imaginary quadratic field the two differ.
\emph{(Proof sketch: $|\mu_K|\in\{2,4,6\}$ with 4 only at $d=1$ and 6 only at
$d=3$; $|\mathrm{disc}(K)|\ge3$ always, $=3$ only at $d=3$, $=4$ only at $d=1$;
case-checking the three unit-group orders leaves $d=1$ as the sole coincidence.)}
The trivial-multiplier criterion is load-bearing: a ``match'' here requires the
natural root to equal the target directly. Under a looser rational-multiplier
criterion, 20 additional non-canonical matches exist in the tested grid, and that
reading is a \etag{sel}, not a theorem. The proof is a complete case analysis over
all imaginary quadratic fields; the $d\le200$ computation is a redundant
cross-check.

\runin{$\Q(\Gs)$ is $\pi$-free}, conditional on Chudnovsky
1976.\margetag{thm} $\Q(\Gs)\cap\Q(\pi)=\Q$: the field generated by $\Gs$ shares
no algebraic content with $\Q(\pi)$ beyond the rationals. The proof reduces a
hypothetical polynomial relation between $\Gs$ and $\pi$ (via the identity
$\Gs\cdot\pi\cdot\sqrt2=\Gamma(1/4)^2$) to a polynomial relation between
$\pi$ and $\Gamma(1/4)$, which Chudnovsky's 1976 algebraic-independence
theorem~\cite{chudnovsky1984} forbids. ``Conditional'' here means the rigour
equals that of this published, foundational transcendence result (consolidated in
Waldschmidt~\cite{waldschmidt2000}), not that it rests on a conjecture.

\runin{Tower-discriminant transcendence}, conditional on
Schneider--Chudnovsky~\cite{schneider1937,chudnovsky1984}.\margetag{thm} For the
tower of the harmonic invariant, the discriminant factors as $\towerdisc$. Then
$A_k$ is rational at $k=3$ ($A_3=1$) and \textbf{transcendental over $\Q$} (i.e.\
$A_k\notin\overline{\Q}$) at every $k\ge4$ ($A_4=4\Gs-1$, etc.): a non-rational
polynomial with rational coefficients in the transcendental $\Gs$ takes
transcendental values.

\runin{BCC complex structure.}\margetag{thm} The 8 BCC corners under $90^\circ$
rotation form two orbits of size~4, and, as a module over the cyclic group
$C_4=\langle 90^\circ\text{ rotation}\rangle\cong\Z/4$,
\[ \bccdecomp, \]
where $V_{\mathrm{complex}}$ carries a natural $\Zi$-module structure
$\cong\Zi^2$. Paired with it is a clean no-go: there is \textbf{no} injective
homomorphism $\Ziu\to O_h^{\mathrm{ab}}$, since $\Ziu\cong\Z/4$ has an order-4
element but $O_h^{\mathrm{ab}}\cong\Z/2\times\Z/2$ (Klein four) does
not, a one-line group-order argument. Together these say the substrate's
$\Zi$ structure lives genuinely in the BCC layer but does \textbf{not} extend to a
global lattice-symmetry embedding. Verified in exact rationals.

\runin{Lemniscatic $L$-value}, conditional on Rubin 1991.\margetag{thm} The
central $L$-value of the lemniscatic curve $E_{\mathrm{lemn}}:y^2=x^3-x$ (Cremona
32.a3) has the clean closed form
\[ \Lvalueeq, \]
via the full BSD formula for the CM rank-0 case~\cite{rubin1991}. (An earlier
$\varpi/2$ came from a BSD-convention double-count; the corrected factor is
$\varpi/4$, verified to 27 digits against the curve's entry in
LMFDB~\cite{lmfdb}.)

\runin{Three further results, conditional on the classical theorems named.}%
\margetag{thm}
The character $\chi_{-4}$ four-level unification\,---\,$\chi_{-4}$ on
$(\Z/4\Z)^\times$ generating the entire $\Gs/\GG$ identity algebra through four
functorial projections (lattice, Chowla--Selberg, Hecke, Dirichlet)\,---\,holds
conditional on Deligne's period conjecture~\cite{deligne1979} in the CM case,
itself proved unconditionally (Blasius/Anderson/Shimura). The $\eta$-tower across
the nine class-number-one imaginary-quadratic fields,
\[ |\eta(\tau_K)|^{2w_K} \;=\; \frac{G_K^{\,w_K}}{(2\pi|d_K|)^{w_K/2}}, \]
holds conditional on Chowla--Selberg~\cite{chowlaSelberg1967}, unifying the
Heegner near-integer phenomenon as a $\chi$-projection. Here $G_K$ is the
field-specific Chowla--Selberg period (it equals $\Gs$ only for the Gaussian
field $d=1$), so the nine fields share the \emph{form} of the identity, not
the constant. The
$\mathrm{Sym}^2\oplus\mathrm{Sym}^3$ structure singles out the construction's
pair. Among the period polynomials $x^2-16\Gpow{a}x+16\Gpow{b}$ with $a<b$
positive integers and positive discriminant (which requires only
$b<2a+\log 4/\log\Gs\approx 2a+1.28$, i.e.\ $b\le 2a+1$ for integers), the pair
$(a,b)=(2,3)$ is distinguished by the \emph{parity of $\Gs$ beneath its
discriminant's radical: $\Delta_{(2,3)}=64\Gsss(4\Gs-1)$ carries an odd power of
$\Gs$}, so $\sqrt{\Delta}=8\Gs\,\sqrtwall$ introduces the generator
$\sqrtwall$ (the very square root that walls Part~II), whereas the
neighbouring pairs carry an even power ($(1,3)$ gives
$\sqrt{\Delta}\propto\sqrt{4-\Gs}$; $(1,2)$ and $(2,4)$ have $b=2a$). Equivalently,
the squarefree radicand of $\Delta$ retains a bare factor of $\Gs$ exactly on the
line $b=2a-1$ (the pairs $(2,3),(3,5),(4,7),\dots$, each adjoining the wall
field $\Q(\Gs)(\sqrtwall)$), of which $(2,3)$ is the minimal member with $a<b$.
That much is elementary algebra; the stronger reading (that $(2,3)$ is the
\emph{unique} natural target) is a \etag{sel}\margetag{sel}, argued from
naturalness, not asserted here as a theorem.

\runin{Phase-G geometric Coulomb.}\margetag{thm} The remaining theorem-grade
node is a lattice fact. On the cubic lattice the static potential between two
manifested sites solves the lattice Poisson equation, and its solution is, at
every finite $L$, exactly the periodic-lattice Poisson Green's function
$G_L$. The emergent inverse-square Coulomb tail is therefore lattice geometry,
carrying no fine-structure content of its own: the lattice coupling reads as a
zero-parameter geometric object, $\alpha_r(r,L)=2\,r\,G_L(r)$, not a free
constant. This is a finite-$L$ identity, established directly, with no continuum
limit required.

\runin{The two honestly-tiered results.} Two of the nine numbered results sit
\emph{below} theorem grade and are stated as such:
\begin{itemize}[leftmargin=1.4em, itemsep=4pt]
  \item \textbf{Coefficient-16}: the value-level identity $16=|\Aut(E)|^2$
        is true, but its structural necessity is a \etag{cnj}, as set out in
        \secref{sec:masterquad}.\margetag{cnj}
  \item \textbf{Phase-J partition-function ultralocality}: a
        \etag[ at $L=2$ only]{thm}.\margetag[ at $L=2$ only]{thm} On a $2^3$
        lattice the classical Euclidean action depends on the state field solely
        through $\sum_i s_i^2$ (the manifested-site count) and is invariant under
        spatial permutation of charge placement; this is proved by explicit
        construction. The general-$L$ extension is \textbf{disconfirmed}: the
        $L=2$ ultralocality is a Nyquist-mode counting degeneracy (every non-zero
        momentum has $\sin(k_i)=0$), not a structural property, and at $L\ge3$ the
        action does depend on spatial placement. The $L=2$ restriction is stated
        plainly and is the spine's honest claim.
\end{itemize}

\runin{The canonical accounting.} That completes the spine: \textbf{nine numbered
results, seven theorem-grade} ($\Gs$ identity, master quadratic, CM uniqueness,
Watson identity, Phase-G geometric Coulomb, harmonic-invariant tower, $\Q(\Gs)$
$\pi$-freeness conditional on Chudnovsky) \textbf{plus two tiered below theorem grade}
(coefficient-16 and Phase-J $L=2$). This is the discrete ontology's firm
ground (its answer to Clause~1), and it stands independent of any physics
interpretation. The match to physics is a separate matter, quarantined in
Part~III.

\section{The construction map}\label{sec:map}
% =============================================================================
% dag.tex  --  The construction DAG (Figure in Part I, section "The construction map")
% -----------------------------------------------------------------------------
% Hand-authored, fully vector, monochrome.  Differentiation is by SHAPE + RULE
% (never colour), mirroring the epistemic-tag tiers of the body:
%   double circle = AXIOM seed   diamond (heavy) = the G* hub
%   solid rounded box = theorem-grade   dashed rounded box = below-theorem
%   small square = framework integer    tiny box = (omitted; ledger anchors live in text)
%   solid arrow = "constructed from"     dashed arrow = "necessity conjectured / weak link"
%
% This figure is hand-set so it stays monochrome, legible, and faithful to the
% nine-result spine accounting (seven theorem-grade, two tiered).
% =============================================================================
\begin{figure}[t!]
  \centering
  \resizebox{\textwidth}{!}{%
  \begin{tikzpicture}[
      >={Stealth[length=4pt]},
      seed/.style   ={circle, double, draw, line width=0.7pt, inner sep=2.2pt,
                      font=\footnotesize, align=center},
      intg/.style   ={rectangle, draw, line width=0.5pt, inner sep=3pt,
                      font=\footnotesize, align=center},
      hub/.style    ={diamond, draw, line width=1.2pt, inner sep=2pt, aspect=1.6,
                      minimum width=1.25cm, minimum height=0.95cm,
                      font=\footnotesize\bfseries, align=center},
      thm/.style    ={rectangle, rounded corners=2pt, draw, line width=0.7pt,
                      inner sep=3pt, font=\footnotesize, align=center, text width=2.3cm},
      tier/.style   ={rectangle, rounded corners=2pt, draw, line width=0.5pt,
                      dash pattern=on 2.4pt off 1.8pt, inner sep=3pt,
                      font=\footnotesize, align=center, text width=2.3cm},
      con/.style    ={->, draw, line width=0.6pt},
      wk/.style     ={->, draw, line width=0.5pt, dash pattern=on 2.4pt off 1.8pt,
                      gray!60!black},
      lbl/.style    ={font=\scriptsize\itshape, fill=white, inner sep=1.2pt},
    ]
    % ---- main trunk (axiom -> integer -> hub -> quadratic -> tower) --------
    \node[seed] (zi)  at (4.0, 0.0)  {$\mathbb{Z}[i]$\\[-1pt]\tiny\scshape axiom};
    \node[intg] (four) at (4.0,-1.9) {$4=|\mathbb{Z}[i]^{\times}|$};
    \node[hub]  (gs)  at (4.0,-3.9)  {$G^{*}$};
    \node[thm]  (mq)  at (4.0,-6.9)  {Master quadratic\\ $x^2-16G^{*2}x+16G^{*3}$\\ \textsc{(t2)}};
    \node[thm]  (tow) at (4.0,-9.4)  {Harmonic\\ invariant tower \textsc{(t8)}};

    % ---- right branches off the hub G* ------------------------------------
    \node[thm]  (wat) at (8.2,-3.1)  {Watson identity\\ $W_3=G^{*2}/2\pi$ \textsc{(t5)}};
    \node[thm]  (qg)  at (8.2,-4.9)  {$\mathbb{Q}(G^{*})$ is\\ $\pi$-free \textsc{(t9)}};
    \node[thm]  (cm)  at (8.2,-6.7)  {CM-curve\\ uniqueness \textsc{(t3)}};
    \node[thm]  (bcc) at (8.2,-8.5)  {BCC complex\\ structure};

    % ---- left / lower: the two below-theorem-grade results (dashed) -------
    \node[tier] (c16) at (0.1,-6.9)  {Coefficient\\ $16=4^2$ \textsc{(t4)}\\ \tiny necessity conj.};
    \node[tier] (pj)  at (0.1,-9.4)  {Phase-J ultra-\\ locality \textsc{(t7)}\\ \tiny $L=2$ only};
    \node[thm]  (pg)  at (8.2,-10.3) {Phase-G geometric\\ Coulomb \textsc{(t6)}};

    % ---- edges -------------------------------------------------------------
    \draw[con] (zi)  -- (four);
    \draw[con] (four) -- (gs);
    \draw[con] (gs)  -- (mq);
    \draw[con] (mq)  -- (tow);
    % the "16 = 4^2" feed from the integer into the quadratic (curved, left)
    \draw[con] (four) to[out=200,in=150] node[lbl,pos=0.62]{$16=4^2$} (mq.west);
    % branches off G* -- the two lowest (bcc, pg) route via an eastward waypoint
    % so they descend on the RIGHT of the central column, clear of the mq box.
    \draw[con] (gs) to[out=18,in=180]  (wat.west);
    \draw[con] (gs) to[out=-8,in=170]  (qg.west);
    \draw[con] (gs) to[out=-30,in=170] (cm.west);
    \draw[con] (gs) to[out=-42,in=120] (6.5,-7.5) to[out=-60,in=170] (bcc.west);
    \draw[con] (gs) to[out=-52,in=120] (6.3,-9.0) to[out=-60,in=160] (pg.west);
    % weak link: 16 = 4^2 necessity is conjectured
    \draw[wk] (mq.west) to[out=180,in=0] (c16.east);
    % Phase-J hangs off the lattice axiom (independent ultralocality fact)
    \draw[wk] (zi.west) to[out=180,in=90] (pj.north);
  \end{tikzpicture}}
  \caption[The construction DAG]{\textbf{The construction DAG.}\enspace Shape
    and rule encode epistemic status: the double-ringed circle is the lone
    \textsc{axiom} ($\mathbb{Z}[i]$), the heavy diamond is the hub ($G^{*}$),
    solid boxes are theorem-grade, and dashed boxes are tiered below theorem
    grade (the coefficient-16 necessity, conjectured; Phase-J, a theorem only at
    $L=2$). The chain runs $i \to \mathbb{Z}[i] \to 4 \to G^{*} \to{}$ master
    quadratic $(16=4^{2}) \to{}$ spine; reading the arrows backward traces every
    theorem-grade node to the one structural choice at the root.}
  \label{fig:dag}
\end{figure}

The spine is not a list but a dependency structure: a directed acyclic graph
rooted in the seed. The axiom reading of the lattice's quarter-turn as $\Zi$
(\secref{sec:seed}) deposits the integer~4; $\Gs$ is constructed four ways from
quarter-integer data (\secref{sec:gstar}); the master quadratic $\masterquad$ is
assembled from $\Gs$ and $16=4^2$ (\secref{sec:masterquad}); and the spine
theorems (harmonic tower, CM uniqueness, $\Q(\Gs)$ $\pi$-freeness, the BCC
complex structure, the $L$-value and its companions) hang off $\Gs$ and the
quadratic in turn (\secref{sec:spine}). Reading the arrows backward is the audit
trail: every theorem-grade node traces, through explicit and checkable steps, to
the one structural choice at the root, and to nothing else.

\part{The Boundary}% =============================================================================
% Part II -- The Boundary  (the climax)
% =============================================================================
\noindent Part~II answers Clause~2 in the only form the present theory can
honestly support: \textbf{where does the construction halt, and what conditional
non-forcing statement can be proved there?} It is the monograph's climax, and
the section where the program's honesty is most directly on the line, because
the quantity at stake is the one whose numerical match ($x_+ =
137.036\ldots$, deferred to Part~III) motivates the whole enterprise. The
temptation to call that match a derivation is exactly what this Part refuses.
What it delivers instead is a \emph{mapped wall}: a route-invariant boundary
around the electromagnetic coupling $\alpha$, built from theorem-grade masonry
around one precisely-located, honestly-marked \etag{opn} gap.

\medskip
\noindent State the shape of the result before its content. The headline\,---\,$\Ftheory$
(the five postulates plus the spine) does not force $\alpha$\,---\,is not a single
theorem. It is a route-invariant boundary assembled from three theorem-grade
components plus one irreducible \etag{opn} kernel; the overall no-go is a
\etag{smc}, and the obstruction it sharpens (the framework's central
foundational obstruction) remains a \etag{fo}.\margetag{fo} The honest
one-line summary is the project's own: \textbf{$\alpha$ is dynamical, not
structural.} It is never ``$\alpha$ is derived,'' and it is never ``$\alpha$ is
proven underivable'' in any absolute sense; \secref{sec:dynamical} states exactly
what it is.

\section{The readout problem}\label{sec:readout}
Part~I built the master quadratic $\masterquad$ as a pure object of the algebraic
spine (\secref{sec:masterquad}): its coefficients are integer multiples of powers
of $\Gs$, its roots $x_+=137.036\ldots$ and $x_-=3.024\ldots$ are an exercise in
the quadratic formula, and the whole of it is a theorem. To turn that polynomial
toward physics, one more thing is needed, and naming it sharply is the entire
content of this Part.

\runin{The readout move.} Read the quadratic as the characteristic polynomial of
a $2\times2$ \emph{readout operator} $W$, an admissible operational
functional on the substrate whose dominant eigenvalue is to be the measured
inverse coupling. For its characteristic polynomial to be $P$, the operator must
have
\[ \assembly. \]
This is the \textbf{operator assembly}. The question that decides whether
$\alpha$ is derived or merely matched is exactly:

\begin{boundarybox}[The readout problem \etag{opn}]
Does $\Ftheory$ (the five postulates plus the algebraic spine plus the
$O_h$ representation theory of the BCC corner module) \emph{force} the
assembly $(\Tr,\Det)=(16\Gss,16\Gsss)$, or is that assembly an \emph{imposed
selection}?
\end{boundarybox}

\noindent This is the readout-structure criterion, and the contract a closure
must satisfy is fixed in advance and deliberately strict: an admissible readout
must be stated \emph{before} any physical value is checked, and a closure fails
immediately if it inserts physical $\alpha$, a CODATA value, $g_c$, or $x_+$, or
uses the framework's transfer matrix $M_N(t)$, which is \emph{defined} to
have the master quadratic as its characteristic polynomial, hence circular and
banned. The missing object is not another fit; it is an operational readout rule.
Part~II determines whether the ontology supplies one.

\runin{Why this is the sharp form of the obstruction.} The central foundational
obstruction of the framework's $\alpha$ program comes down to this one structural
question, because everything \emph{around} it is already settled in Part~I. The
constant $\Gs$ is forced (four ways, \secref{sec:gstar}); the coefficient
$16=|\Ziu|^2$ is a theorem at the value level (\secref{sec:masterquad}); the
polynomial's roots are a theorem. If the \emph{assembly} were forced too, $\alpha$
would be \etag{der}. So the entire weight of the derive-or-match verdict rests on
one hinge: whether $\Ftheory$ produces $W$.

\section{The four FTD-native routes}\label{sec:routes}
A boundary cannot be argued from absence, from merely failing to find an
assembly. The canonical attack instead ran \textbf{four FTD-native routes of
differing character} (later shown, in \secref{sec:routeinv}, to be four dialects
of one underlying obligation, so ``0 of 4'' is one structural fact reached four
ways, not four independent failures). Each was first \emph{force-attempted} (build
the strongest honest forward chain to the assembly, $\Gs$ kept symbolic, all
banned moves excluded) and then \emph{adversarially refuted} by a second pass
hunting for any smuggled step:

\begin{table}[h]
  \centering\small
  \renewcommand{\arraystretch}{1.3}
  \begin{tabularx}{\linewidth}{@{}l >{\RaggedRight\arraybackslash}X@{}}
    \toprule
    \textsc{route} & \textsc{FTD-native principle attempted}\\
    \midrule
    \texttt{jtwist} & the $J$-twisted $\zeta$-regularized determinant operator
      (the candidate clean odd source)\\
    \texttt{bcc} & the BCC body-diagonal transfer / response operator
      (triple-cosine Watson)\\
    \texttt{cm} & the lemniscatic-CM arithmetic of $E:y^2=x^3-x$
      (CM by $\Zi$, $\Aut=\mu_4$)\\
    \texttt{novel} & a forced variational principle / period-ring (Hodge)
      valuation / $K$-theory co-realizability\\
    \bottomrule
  \end{tabularx}
\end{table}

\runin{The verdict: 0 of 4 force the assembly.}\margetag[ no-go]{smc} Every
force-pass self-reported a \texttt{gap}; every refute-pass upheld
\texttt{boundary}; the set of cleanly-forced routes is empty. The honest reading
partitions cleanly into what the ontology \emph{does} force and what it does
\emph{not}:
\begin{itemize}[leftmargin=1.4em, itemsep=4pt]
  \item \textbf{What IS forward-forced} (\etag{der}), agreed by every route with
        nothing smuggled:
    \begin{itemize}[leftmargin=1.2em, itemsep=2pt]
      \item the \textbf{even trace} $16\Gss$: $16=|\mu_4|^2$ for
            $E:y^2=x^3-x$; $\Gss=2\pi\cdot G_{\mathrm{BCC}}(0)$, the Watson BCC
            self-energy;
      \item the \textbf{existence of a clean FTD-native odd source}: the
            $J$-twisted determinant ratio
            $\det_\zeta(D_{3/4})/\det_\zeta(D_{1/4})=\Gamma(1/4)/\Gamma(3/4)=\Gs$
            (degree~1; the $\sqrt{2\pi}$ cancels, so there is no forbidden
            $\sqrt\pi$ prefactor). This is \emph{strictly stronger than
            trace-only}: it means $16\Gsss=16\Gss\cdot\Gs$ is
            \textbf{assemblable}, genuinely lifting the bare-parity no-go.
    \end{itemize}
  \item \textbf{What is NOT forced} (\etag{opn}):
    \begin{itemize}[leftmargin=1.2em, itemsep=2pt]
      \item the \textbf{operator assembly itself}: that the \emph{same}
            $2\times2$ readout carrying $\Tr=16\Gss$ also has $\det=16\Gsss$. For
            a $2\times2$ operator, trace and determinant are \textbf{independent
            invariants}; fixing the trace leaves the determinant free. The odd
            scalar $\Gs$ exists, but nothing forces the \emph{gluing}: neither
            that the two invariants belong to one operator, nor that the odd
            scalar lands in the determinant slot rather than anywhere else. This
            residual is precisely the imposed Vieta target of the
            readout-structure criterion.
    \end{itemize}
\end{itemize}
The differing character of the four channels is what makes ``0/4'' load-bearing
rather than anecdotal; and \secref{sec:routeinv} turns the fact that they
reach one obligation into a theorem. The overall no-go is a \etag{smc}, not a
theorem: the routes are adversarial refutation \emph{attempts}; none constructed
the forcing assembly, and none proved one \emph{cannot} exist. Calling the
boundary a theorem on this evidence would overreach. It is a boundary-mapping
deliverable in the exact sense of Clause~2, at honest conjecture grade.

\runin{A representative route, made concrete.} The boundary shows itself cleanly
in the modular dialect. At the self-dual point $\tau=i$ (the fixed point of the
modular $S$-transformation $\tau\mapsto-1/\tau$), the torus partition function's
modular geometry forces $E_6(i)=0$ and $E_4(i)\propto G^{*4}$, so the available
invariants are \emph{even} powers of the period $\Gs$ ($\Gss$ from the
Green's-function variance, $G^{*4}$ from $E_4$). The required determinant
$16\Gsss$ is an \emph{odd} power, which no theta or modular form on the torus
generates without inserting an external scale. The verdict is exactly
\texttt{UNDERDETERMINED}: this route supplies the even trace and cannot force the
odd determinant slot, one instance of the route-invariant pattern the next
section names.

\section{Route-invariance and the theorem-grade components}\label{sec:routeinv}
The deeper attack asked whether the residual gap could be closed by a sharper
instrument; and in failing to close it, established three genuinely
theorem-grade facts plus the reason all four routes land in the same place. These
are the boundary's \emph{masonry}: each is a theorem, and together they hem the
open kernel of \secref{sec:kbind} into a single, precisely-located gap.

\runin{(a) The flip is ruled out.}\margetag{thm} A \emph{flip} would be a
substrate-native operator that genuinely \emph{forces} $\det=16\Gsss$, flipping
the verdict to ``$\alpha$ derived.'' The only geometric candidate is the
$C_3(\langle111\rangle)$ three-plane $\det_\zeta$ product (a $D_6$-symmetric
object built from three cyclically-permuted planes, each contributing one factor
of $\Gs$). It is \textbf{excluded} from the rank-2 readout in two steps:
\begin{itemize}[leftmargin=1.4em, itemsep=2pt]
  \item a definite complex structure $J$ with $J^2=-I$ (required by the
        trace via the BCC $V_{\mathrm{complex}}=\Zi^2$ structure) needs
        $\mult_O(E)=0$ (the character table of $O$ shows no $O$-symmetric
        2-dimensional subspace of the 8-corner module), which forces $O_h$ to
        break to \textbf{one $C_4$ axis}, so $C_3(\langle111\rangle)\notin\mathrm{Stab}$;
  \item the $D_6$ three-plane product requires
        $C_3(\langle111\rangle)\in\mathrm{Stab}$;
  \item but a $4$-fold face axis and a $3$-fold body-diagonal axis generate the
        whole octahedral rotation group, $\langle C_4, C_3\rangle = O$ (order
        24), so demanding both restores unbroken $O_h$ $\Rightarrow$ no localized
        charge $\Rightarrow$ no $V_{\mathrm{complex}}$ $\Rightarrow$ no readout.
\end{itemize}
The two requirements are mutually exclusive from a single preparation. Every
$2\times2$ that \emph{does} realize $(16\Gss,16\Gsss)$ (e.g.\ the companion
form $\bigl[\begin{smallmatrix}0 & -16\Gsss\\ 1 & 16\Gss\end{smallmatrix}\bigr]$) has
its determinant entry \emph{hand-placed}, which is exactly the banned gluing: a
witness for the cheap argument, not a derivation.

\runin{(b) The equivariant restriction closes its own scope.}\margetag{thm} No
$C_3(\langle111\rangle)$-equivariant rank-2 \emph{restriction} of the three-plane
source carries $(\Tr,\Det)=(16\Gss,16\Gsss)$. On the $C_3$-adapted basis
$\R^3=R\oplus C$, the unique rank-2 $C_3$-invariant subspace is the $C$-plane,
whose $C_3$-commutant is $\{xI+yK\}\cong\C$ (Schur). The master-quadratic roots
are real and distinct ($\mathrm{disc}=64\Gsss(4\Gs-1)>0$), and a
$C_3$-equivariant $C$-plane operator has real-distinct spectrum \emph{iff it is
non-real}, which invokes the ambient scalar $i$, supplied only by breaking
$O_h$ to one $C_4$ axis, whence $\langle C_4,C_3\rangle=O$ and there is no
readout. The chain is therefore \textbf{reality $\Rightarrow$ scalar-$i$
$\Rightarrow$ $C_4$ $\Rightarrow$ $O$}: the \emph{same} wall as the rest of
the program, reached from the restriction side, not an independent conjugacy
fact. The same mechanism holds dimension-free: any $C_3$-equivariant linear
reduction of the three-plane object factors through the $C_3$-fixed diagonal
(where $C_3$ acts as the identity), collapsing the determinant grading
$\Gsss\to G^{*1}$; the rank-2 compressions of $\Gs\cdot I_3$ have determinant
$\Gss$ at most, never $\Gsss$. This closes the restriction's scope; the surviving
non-$C_3$-invariant and det-by-fiat cases are exactly the open kernel.

\runin{(c) The reduction is route-invariant.}\margetag{thm} All four attacks
reduce to the \emph{same} obligation in four dialects, and this is not a
coincidence; it follows from one clean field-theoretic fact:
\begin{quote}\itshape
  $\Q(\Gs)$ is the Galois-fixed field of the master quadratic's $\Z/2$.
\end{quote}
\noindent The quadratic's discriminant is $\Delta=64\Gsss(4\Gs-1)$, so its
splitting field over $\Q(\Gs)$ is $\splitfield$, with Galois group $\Z/2$ swapping
$x_+\leftrightarrow x_-$. Every forward-forced \emph{symmetric} FTD-native
datum (the Watson trace $16\Gss$, the $\det_\zeta$ ratio $\Gs$, the
Chowla--Selberg periods) lives in the $\sigma$-fixed subfield $\Q(\Gs)$ and
is therefore \textbf{provably blind to which root is $1/\alpha$}. The concrete
witness: the family $\det=16\Gss\cdot\Gpow{k}$ for $k=0,1,2,3$ gives dominant
roots $139.05\,/\,137.04\,/\,130.68\,/\,105.76$, \emph{all} of them
$\Ftheory$-consistent, and \textbf{nothing in $\Ftheory$ selects $k=1$.} The
regularized traces $\zeta(-1,1/4)=\zeta(-1,3/4)=1/96$ are rational, so the trace
carries zero $\Gs$; and at rank~2 the regularized determinant is the ordinary
product $x_+\cdot x_-=16\Gsss$, so hitting the assembly is a \emph{chosen}
Vieta target, not a forced one.

\medskip
\noindent These three components are why the boundary is a \emph{wall} and not a
mere absence of progress: the obvious forcing candidate is dead~(a), the
equivariant restriction is closed within its scope~(b), and any remaining route is
provably the \emph{same} route in disguise~(c). What is left is one gap, and
\secref{sec:kbind} states it exactly.

\section{The square-root wall}\label{sec:kbind}
The three theorem-grade components hem the obstruction into a single irreducible
obligation. Stated once, plainly:

\begin{boundarybox}[The binding obligation \etag{opn}]
Prove (or refute) that no substrate-native operator construction can
bind the degree-1, $C_3$-agnostic odd scalar $\Gs$ (the $J$-twisted $\det_\zeta$
ratio) into the determinant slot of the \emph{same} $2\times2$ readout that
carries the definite complex structure $i$, with the exponent fixed at exactly
\textbf{1} by the substrate rather than chosen.
\end{boundarybox}

\noindent The basis-free fact is the primary one: trace and determinant are
independent invariants of a rank-2 operator over $\Q(\Gs)$, so the assembly $W$ is
a free choice. The square-root form below is the cleanest single-coordinate
dialect of that one fact, a basis-representative of a basis-free obstruction,
not a coordinate-bound claim (the determinant's parity is itself basis-dependent:
a rescaling migrates the odd content into the trace). In that dialect: to fix
$\alpha$, the substrate must \textbf{natively realize}
\[ \sqrtwall, \]
the \textbf{squarefree generator} of the quadratic extension
$\Q(\Gs)(\sqrt\Delta)/\Q(\Gs)$; since $\Delta=64\Gsss(4\Gs-1)$ we have
$\sqrt\Delta=8\Gs\cdot\sqrtwall$, the irrational quantity that \emph{first
distinguishes} $x_+$ from $x_-$. Every \emph{symmetric} datum lives in $\Q(\Gs)$
and cannot tell the two roots apart (\secref{sec:routeinv}(c)); the square root is
exactly the ingredient that breaks the $\Z/2$ symmetry and picks a root. This is
the discrete-limit-of-the-square-root wall: the precise place the construction
halts. The substrate forces every \emph{symmetric} ingredient of the quadratic;
it does not, on present evidence, produce the \emph{one antisymmetric
beable} (the square root) that would select the physical branch.

\runin{Why it is \etag{opn}, and not closeable today.} The obligation is a
\textbf{universal negative} over substrate-native operator constructions, and it
is not yet a well-posed proposition, because $\Ftheory$ does not yet pin down the
class it quantifies over. Provisionally, take that class to be the $2\times2$
operators on $V_{\mathrm{complex}}=\Zi^2$ generated by the substrate-native
data (the BCC corner module, the $C_4$-winding preparation, and the
$\det_\zeta$ functor) closed under direct sum and restriction to a stabilizer
(the regularized determinant furnishing the characteristic invariants, not itself
a closure operation). Whether this is the \emph{right} closure, whether a
larger one (tensor products, non-self-adjoint generators) would change the
verdict, is itself part of the open work, not settled here. Until that calculus is fixed precisely, ``no
substrate-native operator binds the odd scalar at exponent~1'' quantifies over an
undefined domain, and one cannot prove a universal over a class one has not
finitely specified. Two senses of \etag{opn} are worth separating: an open
\emph{conjecture} is a well-posed proposition awaiting proof; an open
\emph{desideratum} is a property not yet expressible as a proposition because its
quantifier-domain is unspecified. The binding obligation is currently the latter;
fixing the operator calculus would convert it to the former. \textbf{Crucially,
the conditional theorem of \secref{sec:conditional} is insulated from this}: it
rests only on trace and determinant being independent rank-2 invariants over
$\Q(\Gs)$ (which is well-posed and proved), so the boundary's logic does
not wait on the operator calculus.

\runin{The one non-axiomatic exit is currently shut.} The alternative to
axiomatizing the calculus is to \emph{measure} the binding directly on the
substrate. That measurement has been run, and it returned a \etag{cn}: a null
robust to measurement precision, \textbf{0 macroscopic cluster fissions
across 2000 seeds} in the tested regime, from a protocol that would have detected
fissions above its sensitivity floor and saw none. One should not read
``impossible'' into a finite null; but neither route, the axiomatic (no
calculus) nor the empirical (a null result so far), is presently available, and
this is the single remaining obligation of the entire $\alpha$ program.

\section{The conditional theorem}\label{sec:conditional}
What \emph{is} provable, rigorously and unconditionally, is the conditional. It is
the strongest theorem-grade statement the boundary supports, and the precise sense
in which ``$\alpha$ is not forced'' is earned rather than asserted:

\begin{forcedbox}[The conditional theorem \etag{thm}]
Nothing that $\Ftheory$ forces selects the operator assembly
$(\Tr,\Det)=(16\Gss,16\Gsss)$. Every datum the spine forward-forces is
\emph{symmetric}; it lies in the fixed field $\Q(\Gs)$ of the master
quadratic's $\Z/2$, and is therefore blind to which root is $1/\alpha$. With the
trace $16\Gss$ already forward-forced, only the determinant slot is free: the
assembly is one of an infinite family of determinant values over that fixed
trace, all equally compatible with everything $\Ftheory$ forces. To single it out (equivalently,
to realize a beable in $\Q(\Gs)(\sqrtwall)\setminus\Q(\Gs)$) demands a
substrate-native binding law $W$ beyond $\Ftheory$ that fixes the determinant's
odd-$\Gs$ exponent at exactly \textbf{1} from a single stabilizer.
\end{forcedbox}

\noindent The force of this conditional is carried by the \secref{sec:routeinv}
components, above all the Galois-fixed-field fact, which is what makes
``$\Ftheory$ does not select the assembly'' a theorem and not a restatement of
ignorance: the two sides are exhibited explicitly, as realizations, not merely
asserted.
\begin{itemize}[leftmargin=1.4em, itemsep=3pt]
  \item \textbf{Adopting $W$ is coherent}, the master quadratic is the
        explicit realization: once the assembly is posited it is internally
        consistent.
  \item \textbf{Declining $W$ is equally coherent}: $\det=16G^{*4}$ (dominant
        root $\approx130.68$) and $\det=\Gs$ (dominant root $\approx140.04$) are
        explicit alternative assemblies that everything the spine forces equally
        permits.
\end{itemize}
Because both are coherent, no binding law is implied by what $\Ftheory$
forces; and \emph{that}, not any failed search, is the content of ``the
ontology does not force the assembly.'' $W$ is, in the framework's bookkeeping, a
\textbf{6th-postulate-class input}: added to FTD's axioms \emph{and} shown
substrate-native, it would convert $x_+=1/\alpha$ from a \etag{smc} to
\etag[, modulo $W$]{der}. Whether any FTD-native $W$ exists is exactly the binding
obligation of \secref{sec:kbind}, which remains \etag{opn}. The conditional
itself is theorem-grade; the \emph{unconditional} no-go (``no FTD-native $W$
can ever exist'') is the universal negative, which stays open, and is why the
overall boundary is a \etag{smc}.

\medskip
\noindent This is the exact meaning of the shorthand used elsewhere in the
monograph. ``$\Ftheory$ does not force $\alpha$'' means: no selection follows
from the present five-postulate theory plus the spine; a future $W$ would be a
new binding law, not a consequence already hidden in the current axioms. It does
\emph{not} mean that no possible extension of FTD could ever force the coupling.

\section{Conclusion: \texorpdfstring{$\alpha$}{alpha} is dynamical, not structural}\label{sec:dynamical}
The deliverable of Part~II is a \textbf{mapped wall}. In the project's own
phrasing: \textbf{the discrete ontology forces the \emph{menu}} (the even
trace $16\Gss$ from the Watson integral, and the existence of a clean odd source
$\Gs$ from the $J$-twisted $\zeta$-determinant), \textbf{but it does not force
the \emph{dish}}, the assembly of those ingredients into one readout operator $W$.
So $\alpha$'s value rides on an independent convention, not on the postulates.
The assembly is unforced on present evidence: that contrast is \etag{der};
the reading it licenses, \textbf{$\alpha$ is dynamical, not structural}, is the
\etag{smc} that summarizes it, not a proof that no future binding law exists. The
contrast with $N_c$ is exact and
instructive: $N_c=3$ \emph{is} structural, forced from $O_h$ and topology by
four independent routes with no operator-assembly choice required; the value falls
out of the symmetry. $\alpha$ has no such forcing: with the trace forward-forced,
everything the ontology fixes is still consistent with infinitely many determinant
values over that one trace, the same epistemic status the engine assigns its coupling
$g_c$.

\medskip
\noindent Accordingly, the central obstruction remains a \etag{fo}, and
$x_+=1/\alpha$ remains a \etag{smc}; nothing in this Part promotes either. The
wall is built of theorem-grade masonry (the ruled-out flip, the closed-scope
restriction, the route-invariant reduction, and the conditional theorem) around
one clearly-marked open gap: the square root that would break the root-swapping
symmetry. The surviving exits are explicit and few: a 6th-postulate-class input
$W$ that \emph{forces} the assembly, or a fresh discrete-native measurement (and
the one run so far returned a \etag{cn}).

\medskip
\noindent This is exactly what Clause~2 of the governing goal asks for.
\emph{``Rigorously establish the boundary of what the ontology cannot force''} is
not a consolation prize for failing to derive $\alpha$; it is a deliverable in
its own right. Part~II draws that line where the construction actually halts: at
the discrete limit of a square root. A program that only exhibited the spine of
Part~I would be a curiosity; one that maps this wall, in both directions and at
honest grades, is doing the thing the goal was set to do.

\part{The Bridge}% =============================================================================
% Part III -- The Bridge  (physics quarantine)
% =============================================================================
\noindent Part~III is the \textbf{physics quarantine}. Parts~I and~II did their
work in mathematics, naming the physical target in Part~II only to state precisely
what is \emph{not} forced: an algebraic spine the ontology \emph{forces}
(Clause~1), and a route-invariant boundary showing what it does \emph{not} force
without an additional binding law (Clause~2). This Part introduces the one
ingredient withheld until now (the
\textbf{numerical match to physics}) and states it, and the broader physics
layer around it, at the exact grade each piece carries. It \textbf{derives
nothing}; it states, with the grades foregrounded.

\medskip
\noindent The discipline here is stricter than anywhere else in the monograph,
because physics is where overclaim is easiest and most consequential. The rule for
this Part, stated once and enforced throughout: \textbf{no sentence attaches a
\etag{thm} or a \etag{der} to $\alpha$, to a particle mass, to a mass formula, or
to a gauge ratio.} The $\alpha$ match is a \etag{smc}; the rest of the
construction layer is \etag{par}, with two lepton-anchored mass formulas one tier
higher, stated at that tier and no higher.

\section{The empirical match}\label{sec:match}
Here is the match that motivates the entire enterprise. The master quadratic's
larger root (Part~I, \secref{sec:masterquad}) agrees with the reciprocal of the
electromagnetic fine-structure constant~\cite{codata2022}:

\begin{conjecturebox}[The empirical match \etag{smc}]
\begin{gather*}
  \xplusval, \\
  \alphainveq, \\
  \ppmmatch.
\end{gather*}
\end{conjecturebox}

\noindent A discrete ontology, built only to be internally coherent, throws up a
polynomial one of whose roots lands on $1/\alpha$ to better than two parts in a
million. That is the seed of the whole program.

\medskip
\noindent\textbf{And it is a \emph{match}, not a derivation.}\margetag{smc} The
identification $x_+\leftrightarrow1/\alpha$ is the framework's single load-bearing
physics claim, and it is a \etag{smc}. Part~II established \emph{why} it cannot be
more: the discrete ontology does not force the operator assembly that would read
$x_+$ as the physical coupling, so $\alpha$ is \textbf{dynamical, not structural}
(\secref{sec:dynamical}). The honest sentence is the conjunction of both halves:
\emph{the number matches to $\ppmval$~ppm, and the ontology does not derive that
it should.} $x_+ = 1/\alpha$ has no numerical-uniqueness support.

\section{Retired and closed-negative}\label{sec:closed}
A construction story is incomplete without its dead ends: preserving them is what
stops a closed route from re-emerging as a fresh ``discovery.'' The negative
results below quarantine the physics claim from its discredited neighbours.

\runin{The $\alpha$-derivation routes are all closed.}\margetag{cn} Every attempt
to \emph{derive} $\alpha$ from FTD's dynamics (rather than match it) has
been run and falsified, and each is preserved to prevent re-attempt: \textbf{R1}
(transverse stiffness), \textbf{R2} (source-current normalization), \textbf{R3}
(two-sector response eigenvalue), \textbf{R4} (projected Dirac matter), the
\textbf{Z-factor reading}, \textbf{RG-running}, \textbf{algebraic combinations},
\textbf{$1/\sqrt d$}, \textbf{Langevin-equipartition}, and the \textbf{monomial
scans}, all \etag{cn}. The list is not a record of failure but a map: it is
the empirical content of Part~II's verdict that $\alpha$ is dynamical, traced
route by route. The single non-derivation route that remains live is the
readout/observable-selection class, whose boundary Part~II mapped.

\section{What would close the gap}\label{sec:close}
The gap is not a mystery; it is a precisely-located obligation, and the honest
statement of the research frontier is the statement of what would discharge it.
Part~II named the two exits, and they are the only two on present evidence:
\begin{enumerate}[leftmargin=1.8em, itemsep=4pt]
  \item \textbf{A 6th-postulate-class binding law $W$.} An additional structural
        input (outside the five postulates) that \emph{forces} the
        operator assembly $(\Tr,\Det)=(16\Gss,16\Gsss)$, equivalently that
        natively realizes the antisymmetric beable $\sqrtwall$ which selects the
        physical root (the binding obligation of \secref{sec:kbind}, \etag{opn}).
        Added to FTD's axiom list \emph{and} shown substrate-native, it would
        convert $x_+=1/\alpha$ from a \etag{smc} to \etag[, modulo $W$]{der}.
        Whether any FTD-native $W$ exists is the open kernel of the entire
        $\alpha$ program.
  \item \textbf{A fresh discrete-native measurement.} A measurement of the binding
        directly on the substrate, bypassing imported continuum machinery. This
        exit is currently shut: the one run so far returned a \etag{cn}, a
        null robust to measurement precision, 0 macroscopic cluster fissions
        across 2000 seeds in the tested regime (\secref{sec:kbind}).
\end{enumerate}
These are the live frontier, stated as obligations rather than promises. Neither
exit is presently realized, and that is exactly why $x_+=1/\alpha$ stands where
it does. The monograph claims no imminent closure; it claims to have located the gap
precisely enough that either exit, if it lands, will be recognizable as a closure
rather than another match.

\section{The honest physics scope}\label{sec:scope}
The $\alpha$ match is the framework's sharpest physics claim. The broader
``construction layer'' (the mass spectrum, the gauge ratios, the gravity and
quantum-mechanics-emergence layers) is weaker, and graded as such.

\runin{The construction layer is parametric.}\margetag{par} The bulk of FTD's
Standard Model coverage consists of \textbf{standard physics formulas filled with
FTD constants}: the six quark masses, the ${\sim}42$ meson and ${\sim}48$
baryon masses, the ${\sim}22$ decay rates and widths, the precision-QED
observables ($g-2$, the Lamb shift), and gauge ratios such as $\sinthetaW$ (3.5\%
from the on-shell value) and $\alphastrong$ (0.63\% error). The functional form is borrowed from
QED, chiral perturbation theory, Regge phenomenology, or the renormalization
group, and FTD supplies the integers or couplings. These are \etag{par}
insertions: calibration cross-checks against known physics, not derivations. The
wider mixing-angle sector is moreover competitor-dense: several ratios face
multiple rational competitors at the same tolerance, and $\sin^2\theta_{13}$ is
effectively a mis-prediction, which is exactly why these are cross-checks, not
independent evidence. Of roughly \textbf{162} catalogued Standard Model
quantities, about \textbf{23 are \etag{der}}, \textbf{129 are \etag{par}}, and
\textbf{10 are \etag{imp} or \etag{sel}}. That distribution (a small
genuinely-derived core and a large parametric cross-check layer) is the
honest shape of FTD's physics reach.

\runin{Two lepton-anchored mass formulas sit one tier higher.}\margetag{smc} The
electron mass
\[ \meformula \quad(0.19\%\text{ error}) \]
and the proton--electron ratio
\[ \mpratio \quad(173\text{ ppm error}) \]
are each a \etag{smc}, not merely \etag{par}: they combine FTD's
Moore-neighbourhood integers $\{\Nbase,\Neff,\Nc\}=\{4,13,3\}$ with $\alpha$ in
structurally-motivated combinations, which is harder to dismiss as a bare rational
fit. But the prefactors and exponents are motivated rather than dynamically
derived, so they are stated at that grade and no higher.

\runin{The gravity layer.} The identification of the engine constant $G_N=0.01$
with the physical Newton constant is a \etag{cn}. What the substrate derivation
actually yields is the gravitational fine-structure ratio for one electron,
\[ \alphaGee \quad(0.38\%\text{ match}), \]
and that result carries one flagged interpretive step (the clock hypothesis), so
it is not advanced here as a clean derivation of $G$.

\runin{Dimensionless vs.\ dimensional: where the falsifiable spine lives.} One
structural fact organizes all of the above. FTD's dimensional predictions (any
value in MeV, seconds, or metres) are \textbf{calibration-conditional}: they pass
through exactly two theorem-enforced anchors, $a_{\mathrm{phys}}\equiv\ell_P$
(length) and $K_B=m_e$ (mass), and are no firmer than those declared calibrations
($m_e$ above consumes the $K_B=m_e$ anchor, which is why it is not a free
dimensional prediction). The \textbf{dimensionless} predictions ($\alpha$, the
lepton mass \emph{ratios} $m_\mu/m_e$ and $m_\tau/m_e$, the mixing
angles) require no calibration and constitute the calibration-independent
falsifiable spine. The dimensionless ratios are where FTD stakes a falsifiable
claim; the dimensional values inherit the status of their calibration; and the
broad mass/gauge layer is a \etag{par} cross-check, not independent evidence.


\part*{Coda}
\addcontentsline{toc}{part}{Coda \textmd{\textnormal{--- the map in both directions}}}
\markboth{Coda}{Coda}   % the Coda has no \section; set its running head explicitly
% =============================================================================
% Coda -- The map in both directions
% =============================================================================
\noindent The construction is complete. It is worth saying, in plain language and
for the reader who came to this document from philosophy of physics rather than
from lattice field theory, what the finished map shows, and what it honestly
does not.

\runin{The two findings, stated together.} A finite, discrete ontology\,---\,five
postulates, no completed infinity, no primitive continuum, together with the
single $\Zi$ reading of its order-4 symmetry (\secref{sec:seed})\,---\,\textbf{forces}
a rich theorem-grade algebraic spine: the constant $\Gs=\Gamma(1/4)/\Gamma(3/4)$
(in four representations of one CM period), the master quadratic and its
roots, the arithmetic of the lemniscatic CM curve, the harmonic-invariant tower,
the $\pi$-freeness of $\Q(\Gs)$: \textbf{seven theorem-grade results} in all
(Part~I). And the \emph{same} ontology, in the precise conditional sense of
Part~II, \textbf{does not select} the electromagnetic coupling $\alpha$ without an
additional binding law: the boundary is route-invariant, built from
theorem-grade masonry around one precisely-located open kernel, and its honest
verdict is that $\alpha$ is \emph{dynamical, not structural}. \textbf{Both
halves are the deliverable.} This is the point that is easiest to miss and most
important to hold: the second finding is not the first finding's failure. A
mapped boundary (\emph{this} is what the ontology reaches, and \emph{this} is the
precise place it stops under the present axioms) is a result in its own right, a
different and non-trivial kind of result, one easy to mistake for a failure. A
program that exhibited only the spine would be a curiosity; one that maps the
wall, in both directions and at honest tags, is doing the thing the goal was set
to do.

\runin{The least-wrong self-assessment.} Stated without flattery: FTD is a
\textbf{philosophy-of-mathematics project with a rigorous algebraic core and
suggestive (not derived) physics connections.} The algebraic core is
genuine: the seven theorem-grade results survive skeptical mathematical review,
are independently checkable, and stand independent of any physical
interpretation. The physics connection is real and is honestly weaker than the
mathematics: the $\alpha$ match is a \etag{smc} backed by strong
structural-uniqueness evidence and no derivation chain; the broader mass-and-gauge
layer is \etag{par}. FTD is therefore \textbf{not} a competitor to QED or general
relativity on their empirical terms; it does not out-predict them, and it
does not claim to. Its contribution is of a different kind: a demonstration,
carried out at theorem grade where it can be and at honest conjecture grade where
it cannot, of exactly how far a discrete ontology reaches toward physics on its
own, and exactly where it needs an additional input it does not contain.

\runin{Where it touches the open questions.} The program brushes against two
genuinely open foundational questions, and it touches each at that question's
honest status, claiming no resolution. On \textbf{the measurement problem}, FTD's
quantum-mechanics-emergence layer (the Born rule, collapse-as-projection) is
\etag{sel}/\etag{opn}, not a solution: the load-bearing step ``probability
$=$ normalized energy density'' is unproven, and the framework records it as
unproven. On \textbf{the status of mathematical structure in physics} (whether
the unreasonable effectiveness of mathematics reflects something forced or
something chosen) FTD offers an unusually crisp data point precisely
\emph{because} it draws the line: here is a case where one can exhibit, with
proofs, that a discrete mathematical ontology forces a specific rich structure,
and can exhibit, with theorem-grade conditional machinery, where the same
ontology stops short of a specific physical coupling. The effectiveness is real
on the forced side and absent on the unforced side, and the boundary between them
is mapped rather than asserted. That is the most the project claims, and it
claims it at the tags the claim can bear.

\runin{The two-clause goal, closed.} The value of this construction is not a
derivation of physics from discreteness; that, Part~II shows, the ontology
does not supply. The value is the \textbf{map}: a faithful drawing of the line of
what a finite discrete ontology reaches, \textbf{honest in both
directions}: every theorem proved where the ontology forces, every boundary
marked where it does not, and no tag promoted to make the picture prettier than
it is. The constant at the centre of it all, $\Gs\approx\Gstarapprox$, is the
lemniscatic ratio $\Gamma(1/4)/\Gamma(3/4)$ throughout, never the distinct
lemniscate constant $\varpi\approx\varpiapprox$. That is the construction, and
that is its honest reach.


\part*{Glossary}
\addcontentsline{toc}{part}{Glossary \textmd{\textnormal{--- terms used in this monograph}}}
\markboth{Glossary}{Glossary}
% =============================================================================
% glossary.tex -- plain-language definitions for uncommon terms in the monograph
% =============================================================================
\noindent This glossary gives plain definitions for terms that carry more weight
in this monograph than they usually do in casual prose. It is not a dictionary of
all the mathematics. It is a guide to the vocabulary that makes the argument
readable.

\section*{Core Vocabulary}
\addcontentsline{toc}{section}{Core vocabulary}
\markboth{Glossary}{Core vocabulary}

\begin{termnotes}
  \item[Actual.]
  The part of the ontology that has resolved into a discrete ternary state:
  $s=-1$, $0$, or $+1$. In ordinary language, it is what has happened, as
  opposed to what the substrate is poised to do.

  \item[Axiom.]
  A starting commitment of the framework. An axiom is not derived; it defines
  the model whose consequences are then studied.

  \item[Beable.]
  A quantity treated as something the substrate can actually realize, not merely
  something an outside observer can calculate. In the alpha discussion, the
  missing beable is the branch-selecting square root.

  \item[Binding law.]
  A rule that would force separate ingredients into one operational readout. For
  $\alpha$, it would have to bind the trace and determinant into the same
  rank-2 operator for an internal reason, rather than by choice.

  \item[Boundary.]
  A place where the construction stops. A boundary is not a failure by itself;
  it is useful information about what the axioms do and do not determine.

  \item[Calibration.]
  A declared bridge between dimensionless lattice units and physical units such
  as metres, seconds, or masses. Once a calibration is accepted, many dimensional
  quantities inherit it.

  \item[Closed negative.]
  A route that has been tried and ruled out within its stated scope. The result
  is preserved so the same failed route does not keep returning as if new.

  \item[Construction.]
  A step-by-step building of one object from earlier objects, like building
  integers from natural numbers or complex numbers from real pairs. In this
  monograph, construction means the reader should be able to audit how each
  object enters.

  \item[Dispositional.]
  The part of the ontology that represents tendency, capacity, or readiness to
  act. Here the flux field is dispositional: it stores what the substrate is
  poised to manifest.

  \item[Dynamical.]
  Determined by an additional rule, process, state, or measurement, rather than
  fixed by the bare structure. Saying $\alpha$ is dynamical means the current
  structural spine does not by itself select its value.

  \item[Epistemic tag.]
  A label that states what kind of claim is being made: theorem, conjecture,
  parametric insertion, open problem, and so on. The tag is part of the claim,
  not a side comment.

  \item[Foundational obstruction.]
  A named obstacle that blocks a hoped-for derivation at the level of principles,
  not merely at the level of algebraic tidying.

  \item[Parametric insertion.]
  A calculation where the formula comes from standard physics and FTD supplies
  numbers to put into it. It can be a useful cross-check, but it is not a
  derivation from the FTD axioms.

  \item[Readout.]
  The rule that turns a mathematical structure inside the framework into an
  operational quantity that could be compared with physics.

  \item[Selection.]
  A motivated choice among possibilities. It may be natural, but unless the
  framework forces it uniquely, it is not a theorem.

  \item[Structural.]
  Fixed by the architecture of the framework itself. A structural value is not
  chosen by measurement or calibration after the fact.

  \item[Substrate.]
  The underlying discrete lattice system: sites, states, local neighborhoods,
  and update rules before any continuum or physical interpretation is added.

  \item[Substrate-native.]
  Built from the substrate's own objects and rules. A substrate-native readout
  should not import the desired physical answer from outside.
\end{termnotes}

\section*{Mathematical Vocabulary}
\addcontentsline{toc}{section}{Mathematical vocabulary}
\markboth{Glossary}{Mathematical vocabulary}

\begin{termnotes}
  \item[BCC layer.]
  The body-centred-cubic layer inside the Moore neighbourhood: the eight body
  diagonal corners around a site. It is important because it carries a natural
  quarter-turn complex structure.

  \item[Branch selection.]
  Choosing one root of a quadratic rather than the other. The symmetric data of
  a quadratic know both roots together; branch selection needs the ingredient
  that distinguishes them.

  \item[CM period.]
  A special number arising from an elliptic curve with complex multiplication.
  In this monograph, $G^*$ is treated as a lemniscatic CM period appearing in
  several different forms.

  \item[Determinant.]
  For a rank-2 readout operator, the product-like invariant. In the alpha
  problem, the determinant slot is exactly where the unforced odd factor of
  $G^*$ would have to land.

  \item[Gaussian integers.]
  Numbers of the form $a+bi$ with $a$ and $b$ integers. They enter because a
  quarter-turn in a coordinate plane behaves like multiplication by $i$.

  \item[Galois-fixed.]
  Unchanged by the symmetry that swaps the two roots of the master quadratic.
  Data that are Galois-fixed can describe the pair of roots but cannot choose
  which one is the physical branch.

  \item[Lemniscatic ratio.]
  The number $G^*=\Gamma(1/4)/\Gamma(3/4)$. It is central to the construction
  and should not be confused with the related classical lemniscate constant
  $\varpi$.

  \item[Master quadratic.]
  The polynomial
  \[
    x^2-16G^{*2}x+16G^{*3}.
  \]
  Its larger root is the number that matches $\alpha^{-1}$; the polynomial
  itself is algebraic, while the physical identification remains conjectural.

  \item[Moore neighbourhood.]
  The 26 neighbours around a cubic-lattice site: face, edge, and corner
  neighbours together. It is the local causal neighbourhood of the model.

  \item[Rank-2 operator.]
  A two-dimensional linear readout object. Here it matters because such an
  operator is controlled by two invariants: trace and determinant.

  \item[Route-invariant.]
  Reached by different lines of attack in the same form. The alpha boundary is
  route-invariant because several natural attempts all reduce to the same
  missing binding obligation.

  \item[Square-root wall.]
  The obstruction represented by $\sqrt{G^*(4G^*-1)}$, the first object that
  distinguishes the two roots of the master quadratic. The current substrate
  does not force this branch-selecting object.

  \item[Trace.]
  For a rank-2 readout operator, the sum-like invariant. In the alpha problem,
  the trace is much better supported than the determinant assembly.

  \item[Zeta-regularized determinant.]
  A way to assign a determinant-like value to an infinite spectrum using zeta
  regularization. In this monograph it supplies a clean source of one factor of
  $G^*$, but not the full readout assembly.
\end{termnotes}

\section*{Phrases Used in the Argument}
\addcontentsline{toc}{section}{Phrases used in the argument}
\markboth{Glossary}{Phrases used in the argument}

\begin{termnotes}
  \item[The menu but not the dish.]
  A shorthand for the alpha boundary. The framework supplies ingredients that
  make the alpha readout plausible, but it does not force their final assembly.

  \item[Physics quarantine.]
  The separation of the mathematical construction from the physics-facing match.
  The phrase is meant to keep a good numerical match from quietly turning into a
  claimed derivation.

  \item[The spine.]
  The chain of theorem-grade algebraic results that can be studied without any
  physical interpretation.

  \item[The wall.]
  A memorable name for a mapped obstruction. It says: this is where the current
  construction stops, and this is the exact form of what is missing.
\end{termnotes}


\part*{Reader's Notes}
\addcontentsline{toc}{part}{Reader's notes \textmd{\textnormal{--- reading the boundary}}}
\markboth{Reader's Notes}{Reader's Notes}
% =============================================================================
% reader_notes.tex -- reader-facing notes, no project-path references
% =============================================================================
\noindent The body has deliberately moved quickly through a lot of material:
axioms, algebraic constructions, a numerical match, and then a boundary where a
tempting derivation stops short. These notes slow down at the places a careful
reader is most likely to ask, ``Wait, what exactly is being claimed here?'' They
do not add new results; they unpack the rules of reading the results already
given.

\section*{How to read the tags}
\addcontentsline{toc}{section}{How to read the tags}
\markboth{Reader's Notes}{How to read the tags}

\noindent The epistemic tags are not decorative labels. They are the grammar of
the argument. A \etag{thm} tag says: if you accept the axioms and the cited
classical mathematics, the statement follows. A \etag{par} tag says: a standard
physics formula has been supplied with FTD inputs, so the numerical agreement is
a cross-check rather than a derivation. A \etag{smc} tag says: the evidence is
substantial enough to matter, but there is still no chain from the axioms to the
physical identification. An \etag{opn} tag says: the question is live, and the
work needed to settle it has been named.

\medskip
\noindent This is why the same page can say that a polynomial is theorem-grade
while its physical interpretation is not. The algebraic statement
``this quadratic has these roots'' is a different species of claim from the
physics statement ``the larger root is the measured inverse electromagnetic
coupling.'' The first is a calculation inside the construction. The second is a
bridge to nature, and it has to earn a stronger readout rule before it can be
called a derivation.

\section*{Why the alpha match is not a derivation}
\addcontentsline{toc}{section}{Why the alpha match is not a derivation}
\markboth{Reader's Notes}{Why the alpha match is not a derivation}

\noindent The large root $x_+$ lands very close to $\alpha^{-1}$, and the match
is good enough that dismissing it casually would be intellectually lazy. But a
match, even a striking one, is not the same thing as a mechanism. To derive the
fine-structure constant, the framework would need to show why the physical
electromagnetic readout must use the larger root of the master quadratic.

\medskip
\noindent The missing step is not a better numerical approximation. The missing
step is an operational rule. In the language of Part~II, the framework must
produce a rank-2 readout operator whose trace and determinant are exactly
\[
  (\Tr,\Det)=(16G^{*2},16G^{*3})
\]
for a structural reason internal to the substrate. The trace and the odd
$G^*$ factor both have native sources. What is not forced is their assembly into
the same operator, with the odd factor placed specifically in the determinant
slot. That placement is the difference between a discovered readout and a chosen
one.

\section*{What the square-root wall means}
\addcontentsline{toc}{section}{What the square-root wall means}
\markboth{Reader's Notes}{What the square-root wall means}

\noindent A quadratic has two roots, and the symmetric data of the quadratic do
not by themselves choose between them. The trace and determinant know the pair of
roots together; they do not tell us which root nature should read as
$1/\alpha$. The quantity
\[
  \sqrt{G^*(4G^*-1)}
\]
is the first ingredient that distinguishes the two branches. That is why the
boundary is described as a square-root wall: the framework can build the
symmetric side of the quadratic, but it does not yet natively produce the
antisymmetric branch-selection object.

\medskip
\noindent Put less technically: FTD can get to the menu of possible readings,
but it does not force the dish. The present theory allows the alpha-like reading
and also allows other determinant choices over the same trace. Because both
sorts of readout are coherent, the alpha reading is not forced by the current
postulates alone.

\section*{What would close the gap}
\addcontentsline{toc}{section}{What would close the gap}
\markboth{Reader's Notes}{What would close the gap}

\noindent There are two honest ways the gap could close. The first would be a
new binding law: a substrate-native rule that fixes the readout assembly rather
than choosing it after seeing the target. If such a law were added and justified,
the alpha match would become a derivation relative to that added law.

\medskip
\noindent The second would be a direct substrate measurement of the binding
phenomenon. That route has not produced a positive result in the tested regime.
A finite null is not a universal impossibility theorem, but it does mean the
empirical path has not supplied the missing readout either.

\medskip
\noindent So the present conclusion is deliberately modest. The monograph does
not say that no future extension could ever derive $\alpha$. It says that the
current five-postulate ontology plus the algebraic spine does not already contain
the binding law that would make the derivation go through.

\section*{How to read the physics layer}
\addcontentsline{toc}{section}{How to read the physics layer}
\markboth{Reader's Notes}{How to read the physics layer}

\noindent The broader physics catalogue should be read with the same discipline.
Dimensionless matches are the most informative, because they do not depend on a
choice of metres, kilograms, seconds, or calibration anchors. Dimensional
quantities are weaker: once a length scale or mass scale is declared, many
numbers inherit that declaration.

\medskip
\noindent Likewise, a particle mass or decay rate computed by inserting FTD
numbers into a familiar formula from quantum field theory is not automatically an
FTD derivation. It may still be useful. It may show that the framework's
integers and constants sit in the right numerical neighborhood. But the
functional form is borrowed, so the honest tag is \etag{par}. The point of the
tagging system is not to diminish those matches; it is to keep them in the
right drawer.


\backmatter
\printbibliography[heading=bibintoc, title={Canonical sources \& references}]

% ---- colophon -------------------------------------------------------------
\clearpage
\thispagestyle{empty}
\vspace*{\fill}
\begin{center}\small\itshape
Colophon
\end{center}
\begin{center}\footnotesize
This monograph was set in Times-like type with \LaTeX. Its design is
deliberately monochrome: epistemic tier is carried by rule weight and style
(heavy, thin, dashed, double), never by colour, so the argument reads
identically in grayscale and in print, after Bringhurst's counsel that structure
should be legible without ornament~\cite{bringhurst2004}. The typesetting is
meant to support the same discipline as the prose: every result is shown at the
grade it earned, and none is made stronger by presentation.
\end{center}
\vspace*{\fill}

\end{document}
